\begin{textsample}{POS Dim 4 – 2018\_c6 – Score 12.00 – 2018\_c6/enade\_2018\_publi\_6e03840cda.txt}  \label{ex:f4_pos_012}
CAPÍTULO 6 ANÁLISE TÉCNICA DA PROVA Este capítulo tem por objetivo apresentar o desempenho dos estudantes concluintes de Comunicação Social – Publicidade e Propaganda no Enade/2018.
Para isso, foram calculadas as estatísticas básicas da prova, em seu todo ( seção 6.1.1 ), bem como as estatísticas relacionadas aos Componentes de Formação Geral ( seção 6.1.2 ) e de Conhecimento Específico da Área ( seção 6.1.3 ).
Dadas as suas características, foram analisadas, em separado, as questões objetivas ( seção 6.2 ) e as questões discursivas ( seção 6.3 ).
Para as questões objetivas, foram disponibilizados os índices de facilidade e de discriminação Ponto-Bisserial, também em separado, para os Componentes de Formação Geral ( seção 6.2.1 ) e os de Conhecimento Específico ( seção 6.2.2 ).
De cada componente, uma das questões foi escolhida para exemplificar a análise gráfica, relacionando -se as alternativas escolhidas pelos estudantes ( inclusive o gabarito ) com o número de acertos no componente.
No Anexo I, apresenta -se a íntegra da análise gráfica para todas as questões objetivas.
Para cada uma das questões discursivas, os conteúdos dos tipos mais comuns de respostas dos estudantes são apresentados e comparados com o padrão de respostas esperado ( ver Anexo VIII com o padrão de respostas ).
Tomando -se como base as duas questões discursivas do Componente de Formação Geral, na seção 6.3.1.6, são apresentados comentários sobre a correção das respostas em relação à Língua Portuguesa.
Nas tabelas, constam as seguintes estatísticas das notas²⁰ : média do desempenho na prova, erro padrão da média, desvio padrão, nota mínima, mediana e nota máxima para cada um de seus componentes.
Tais estatísticas contemplam o total de estudantes concluintes da Área de Comunicação Social – Publicidade e Propaganda inscritos e presentes à prova do Enade/2018, tendo em vista agregações, ou por Grandes Regiões e o país como um todo, ou por Categoria Administrativa, Organização Acadêmica da IES e Modalidade de Ensino.
Em relação aos gráficos de distribuição de notas, o intervalo considerado foi de 10 unidades, aberto à esquerda e fechado à direita, com exceção do primeiro intervalo, [ 0 ; 10 ], fechado em ambos os extremos.
Para os gráficos de distribuição das notas das questões discursivas, foram consideradas mais duas categorias : questão em branco²¹ e nota zero. ²⁰ Uma definição dessas estatísticas pode ser encontrada no Glossário. ²¹ Nesse grupo estão incluídas também as respostas classificadas como nulas ou desconsideradas. 6.1 ESTATÍSTICAS BÁSICAS DA PROVA Nesta seção, apresentam -se estatísticas selecionadas e histogramas da nota geral ( 6.1.1 ) e de cada componente : Formação Geral ( 6.1.2 ) e Conhecimento Específico ( 6.1.3 ).
São também apresentadas estatísticas selecionadas de subpopulações, caracterizadas por Grande Região, Categoria Administrativa, Organização Acadêmica e Modalidade de Ensino. 6.1.1 Estatísticas Básicas Gerais Na Tabela 6.1, são apresentadas as Estatísticas Básicas da prova, por Grande Região, dos estudantes concluintes de Comunicação Social – Publicidade e Propaganda.
Observa -se que não existem cursos dessa Área em CEFET/IFET, nem na modalidade Educação a Distância.
A Média das notas da prova, como um todo, nas seções seguintes serão analisados os Componentes de Formação Geral e de Conhecimento Específico foi 50,0, sendo que os alunos da região Norte obtiveram a Média mais baixa ( 48,0 ), e os da região Sul obtiveram a Média mais alta ( 52,9 ).
As demais médias foram : 48,5 na região Nordeste, 49,6 na região Sudeste e 50,9 na região Centro-Oeste.
O Desvio padrão para o Brasil, como um todo, foi 14,2, sendo o maior Desvio padrão encontrado na região Centro-Oeste ( 14,6 ), e o menor, nas regiões Norte e Nordeste ( 13,8 ), indicando uma dispersão um pouco menor das notas dessas últimas regiões.
A região que obteve a maior nota Máxima foi a Sul ( 89,8 ), ao passo que a região que atingiu a menor nota Máxima foi a Norte ( 78,4 ).
A Mediana do Brasil, como um todo, foi 50,8, sendo a maior Mediana obtida na região Sul ( 54,2 ), e a menor, obtida na Norte ( 48,0 ).
A nota Mínima do Brasil, zero, foi obtida nas regiões Sudeste, Sul e Centro-Oeste.
A Nota Mínima na região Nordeste foi 6,9, e a maior nota Mínima entre as cinco Grandes Regiões, foi 16,8, na região Norte.
Considerando -se as notas, segundo Grande Região, observa -se que existe diferença estatisticamente significativa ao nível de 95 \% entre a maior Média, obtida na região Sul, e as demais.
As diferenças entre a Média da região Centro-Oeste e as médias das regiões Norte e Nordeste, bem como a diferença entre as médias das regiões Sudeste e Nordeste são estatisticamente significativas22. 22 Todas as comparações deste capítulo utilizam os intervalos de 95 \%.
Os erros-padrão da média que possibilitam os testes estão disponíveis nas tabelas.
O comportamento das notas dos estudantes de todo o Brasil pode ser observado no Gráfico 6.1, que apresenta um histograma com a distribuição das mesmas : uma distribuição unimodal com a moda no intervalo ( 50 ; 60 ], seguido de perto pelo intervalo ( 40 ; 50 ].
Na Tabela 6.2, são apresentadas informações referentes à Média da nota final, desagregadas por Categoria Administrativa, por Organização Acadêmica e por Modalidade de Ensino.
A Média das notas da prova, como um todo, foi 50,0.
Em relação à Categoria Administrativa, os estudantes das IES Públicas obtiveram Média mais alta ( 56,2 ) que a Média nacional, e os das IES Privadas obtiveram Média mais baixa ( 49,5 ).
Observa -se que existe diferença estatisticamente significativa entre as médias das notas das IES Públicas e as das Privadas.
A diferença entre as médias das regiões Sul e Norte ( 4,9 ), a maior e a menor Média, é inferior à diferença entre a das IES Públicas e a das Privadas ( 6,7 ), caracterizando -se maior diversidade administrativa do que regional.
O Desvio padrão para as IES Públicas ( 14,2 ) foi igual ao do Brasil, como um todo, e o das IES Privadas foi ligeiramente menor ( 14,1 ).
No tocante à Organização Acadêmica, as Universidades e as Faculdades obtiveram Média mais alta que a nacional ( 50,8 e 50,3, respectivamente ).
A Média dos Centros Universitários foi menor do que a nacional ( 47,9 ).
Constata -se que existe diferença estatisticamente significativa ao nível de 95 \% entre a Média mais baixa, 47,9, dos Centros Universitários, e a dos demais tipos de Organização Acadêmica.
A Organização Acadêmica CEFET/IFET não teve alunos inscritos e presentes.
A modalidade de Ensino Educação a Distância não teve estudantes realizando o exame.
Sendo assim, não haverá a análise por modalidade. 6.1.2 Estatísticas Básicas no Componente de Formação Geral Na Tabela 6.3, são apresentadas as Estatísticas Básicas em relação ao componente da prova que avalia a Formação Geral dos estudantes concluintes.
Os alunos de todo o Brasil obtiveram desempenho médio de 48,2.
Quanto à variabilidade, o Desvio padrão das notas dos estudantes do Brasil, como um todo, foi 17,4.
A maior Média foi obtida na região Sul ( 50,0 ), e a menor, na região Nordeste ( 46,3 ).
As demais médias foram : 46,6 na região Norte, 48,1 na região Sudeste e 49,9 na região Centro-Oeste.
Já o maior Desvio padrão foi obtido na região Centro-Oeste ( 18,0 ), e o menor, na região Norte ( 16,8 ).
Os demais desvios padrões foram : 17,0 na região Nordeste, 17,4 na região Sudeste e 17,5 na região Sul.
A maior nota no Componente de Formação Geral da prova do Enade/2018 foi 96,8, obtida por, pelo menos, um aluno na região Sudeste.
A menor nota Máxima foi obtida na região Norte ( 86,8 ).
A Mediana do Brasil, como um todo, foi 47,8, sendo a menor Mediana encontrada na região Norte ( 45,2 ), e a maior, na região Centro-Oeste ( 50,4 ).
A nota Mínima nessa parte foi zero em quase todas as regiões, com exceção da região Norte cuja nota Mínima foi 4,6.
Considerando -se as notas, segundo Grande Região, observa -se que existe diferença estatisticamente significativa entre as Médias das notas do Componente de Formação Geral obtidas nas diversas regiões.
Não são estatisticamente significativas as diferenças entre a Média da região Norte e a das regiões Nordeste e Sudeste, e entre a Média da região Centro-Oeste em relação às das regiões Sudeste e Sul.
No Gráfico 6.2, pode -se verificar a avaliação do desempenho dos estudantes do Componente de Formação Geral, a partir do histograma da distribuição das notas correspondentes.
A distribuição é unimodal, com moda em ( 40 ; 50 ], intervalo imediatamente anterior ao modal da distribuição de notas da prova, como um todo ( Gráfico 6.1 ).
Nota -se, ainda, que, no Gráfico 6.2, as notas apresentam maior dispersão do que as no Gráfico 6.1 ( distribuição das notas da prova ), confirmado pela comparação dos desvios padrões : 14,2 para a nota da prova, como um todo, e 17,4 para o Componente de Formação Geral.
Na Tabela 6.4, são apresentadas as informações referentes ao desempenho dos concluintes do Componente de Formação Geral, em diferentes agregações : Categoria Administrativa e Organização Acadêmica.
Observa -se que existe diferença estatisticamente significativa entre as médias dos tipos de Categoria Administrativa : a maior Média obtida por estudantes de IES Públicas ( 56,5 ), e a menor, pelos de IES Privadas ( 47,5 ).
Nota -se que também há diferença estatisticamente significativa entre a maior das médias dos três tipos de Organização Acadêmica e a menor delas sendo, o valor maior para as Universidades ( 49,2 ) e o menor para os Centros Universitários ( 45,4 ). 6.1.3 - Estatísticas Básicas do Componente de Conhecimento Específico Na Tabela 6.5, são apresentadas as Estatísticas Básicas referentes ao Componente de Conhecimento Específico da área de Comunicação Social – Publicidade e Propaganda.
A Média do desempenho dos alunos do Brasil, como um todo, foi 50,6.
A maior Média foi obtida na região Sul ( 53,9 ), e a menor, na região Norte ( 48,4 ).
As demais médias foram : 49,3 na região Nordeste, 50,1 na região Sudeste e 51,2 na região Centro-Oeste.
Quanto à variabilidade das notas, o Desvio padrão do Brasil, como um todo, foi 15,1, sendo o maior Desvio padrão observado na região Centro-Oeste ( 15,3 ), e o menor, na região Nordeste ( 14,6 ).
Os demais desvios foram : 14,9 na região Norte, 15,2 na região Sudeste e 14,8 na região Sul.
A Mediana das notas dos estudantes de todo o Brasil foi 52,0.
A maior Mediana ocorreu na região Sul ( 55,5 ), e a menor, na região Norte ( 49,0 ).
As demais medianas foram : 50,0 na região Nordeste, 51,5 na região Sudeste e 52,6 na região Centro-Oeste.
A nota Máxima do Brasil, como um todo, foi 90,0, sendo obtida por, pelo menos, um aluno na região Nordeste.
As demais notas máximas foram : 78,8 na região Norte, 88,5 na região Sudeste, 89,9 na região Sul e 84,9 na região Centro-Oeste.
A nota Mínima foi zero nas regiões Sudeste, Sul e Centro-Oeste.
As notas Mínimas nas regiões Norte e Nordeste foram, respectivamente, 12,2 e 2,5.
Observa -se que não existe diferença estatisticamente significativa entre as duas menores Médias das notas do Componente de Conhecimento Específico, nas regiões Norte e Nordeste.
No entanto, existe diferença estatisticamente significativa entre a região Sul, que tem a maior média, e as demais regiões.
Assim como nos Gráficos 6.1 e 6.2, no Gráfico 6.3, apresentado a seguir, é possível ser feita uma avaliação do desempenho de concluintes de Comunicação Social – Publicidade e Propaganda em relação ao Componente de Conhecimento Específico, com um histograma da distribuição das notas correspondentes.
Esta também é uma distribuição unimodal, e o grupo modal é o ( 50 ; 60 ], igual ao grupo modal da prova, como um todo, e acima do grupo modal para a Formação Geral.
Na Tabela 6.6, é apresentada uma comparação dos resultados em relação à Categoria Administrativa e à Organização Acadêmica, agora levando -se em conta o desempenho de alunos do Componente de Conhecimento Específico da prova.
No que se refere à Organização Acadêmica, a maior Média foi a das Universidades ( 51,3 ), vindo a seguir a das Faculdades ( 50,9 ) e, por fim, a dos Centros Universitários ( 48,7 ).
O Desvio padrão, 15,1, foi o mesmo para os três tipos de Organização Acadêmica e igual ao valor para o Brasil, como um todo.
As Faculdades obtiveram a maior nota Máxima ( 90,0 ).
Os Centros Universitários obtiveram nota Máxima 89,9, enquanto as Universidades obtiveram nota Máxima 86,9.
As medianas foram : 52,6 nas Universidades, a maior delas ; 49,5 nos Centros Universitários, a menor delas ; e 52,5 nas Faculdades.
A nota Mínima foi zero para todas as Organizações Acadêmicas.
Observa -se que existe diferença estatisticamente significativa ao nível de 95,0 \% no Componente de Conhecimento Específico entre as notas dos Centros Universitários e as notas tanto das Universidades quanto das Faculdades.
Quanto à Categoria Administrativa, observa -se um comportamento semelhante àquele da parte de Formação Geral e à prova, como um todo, ou seja, existe diferença estatisticamente significativa entre as médias das IES Públicas ( 56,1 ) e as das IES Privadas ( 50,1 ).
Nesse caso, também, a maior Média foi obtida por alunos de IES Públicas de ensino. 6.2 ANÁLISE DAS QUESTÕES OBJETIVAS Esta seção apresenta estatísticas selecionadas e histogramas das Questões Objetivas de Formação Geral ( 6.2.1 ) e de Conhecimento Específico ( 6.2.2 ).
São também apresentadas e comparadas as médias das subpopulações, caracterizadas por Grande Região. 6.2.1 Componente de Formação Geral Na Tabela 6.7, são apresentadas as Estatísticas Básicas relativas às oito questões objetivas do componente da prova que abrange a Formação Geral dos estudantes.
A Média do Brasil foi 54,9.
A menor Média foi encontrada na região Nordeste ( 52,8 ), e a maior, na região Centro-Oeste ( 58,3 ).
As demais médias foram : 53,6 na região Norte, 54,5 na região Sudeste e 57,4 na região Sul.
O Desvio padrão do Brasil foi 23,6, sendo o maior Desvio padrão encontrado na região Sudeste ( 23,8 ), e o menor, na região Nordeste ( 22,9 ).
Os demais desvios foram 23,7 na região Norte e 23,4 nas regiões Sul e Centro-Oeste.
As medianas do Brasil, como um todo, e de três das cinco regiões foram 50,0.
As regiões Sul e Centro-Oeste apresentaram Mediana igual a 62,5.
A nota Máxima foi 100,0 em todas as regiões.
A nota Mínima ( 0,0 ) foi igual para todas as regiões.
Na Tabela 6.8, são apresentados o Índice de Facilidade e o Índice de Discriminação ( Ponto-Bisserial ) para cada uma das questões objetivas do Componente de Formação Geral.
Quanto ao Índice de Facilidade, foram usadas as seguintes cores para diferenciar o nível de dificuldade da questão :  Azul para as questões classificadas com índice Muito fácil ( \&gt ; =0,86 ), verde para as questões classificadas com índice Fácil ( 0,61 a 0,85 ), amarelo para as questões classificadas com Médio ( 0,41 a 0,60 ), vermelho para as questões classificadas com Difícil ( 0,16 a 0,40 ) e roxo para as questões classificadas com Muito difícil ( \&lt ; =0,15 ).
Já quanto ao Índice de Discriminação, foram usadas as seguintes cores para qualificar a questão :  As questões classificadas com índice Fraco receberam a cor vermelha ( \&lt ; =0,19 ), as classificadas com Médio receberam a cor amarela ( 0,20 a 0,29 ), as classificadas com Bom receberam a cor verde ( 0,30 a 0,39 ) e as classificadas com Muito bom ( \&gt ; =0,40 ) receberam a cor azul.
As questões objetivas do Componente de Formação Geral, segundo o Índice de Facilidade, foram assim avaliadas : das oito questões, nenhuma teve o Índice de Facilidade classificado como Muito fácil, e duas questão foram tidas como Fácil, por terem índice de acertos 0,61 e 0,85.
Cinco questões foram consideradas com índice de dificuldade Médio, situando -se no intervalo entre 0,41 e 0,60 do Índice de Facilidade, ou seja, houve entre 41,0 \% e 60,0 \% de acertos, enquanto uma questão foi classificada na categoria Difícil, com Índice de Facilidade 0,36.
Nenhuma questão apresentou menos de 15 \% de acertos, razão pela qual não houve questão classificada como Muito difícil.
O Índice de Facilidade variou de 0,36 a 0,67, e o de Discriminação, de 0,44 a 0,53.
Na Tabela 6.9, é apresentada a distribuição das questões, levando -se em conta, simultaneamente, a classificação dos dois índices.
Todas as oito questões tiveram Índice de Discriminação Muito bom e figuraram entre três níveis de dificuldade Fácil, Médio e Difícil ; sendo duas classificadas na categoria Fácil ( questões 2 e 8 ) do Índice de Facilidade, cinco na categoria Médio ( questões 1, 3, 5, 6 e 7 ), e uma, na categoria Difícil ( a questão 4 ).
Em particular, a questão 4 foi a que apresentou o maior poder discriminatório, com índice 0,53, e foi considerada Difícil em termos de facilidade, com uma proporção de 0,36 acertos.
O máximo de acertos foi alcançado pela questão 8, com um Índice de Facilidade de 0,67, seguida pela questão 2, com um Índice de Facilidade de 0,66.
Já os índices de Discriminação foram Muito bom para as todas as questões.
No Gráfico 6.4, para exemplificar, analisa -se o comportamento da questão de número 4 de Formação Geral.
Trata -se de uma questão considerada Difícil, em relação à facilidade e a que obteve o maior Índice de Discriminação dessa parte da prova.
Neste gráfico, cada uma das cinco curvas representa o percentual de respostas em determinada alternativa da questão, em função do número de acertos dos estudantes nessa parte da prova ( Formação Geral/Múltipla Escolha ), antes de possíveis eliminações pelo critério do Ponto-Bisserial.
A curva em verde corresponde à alternativa B, a alternativa correta para essa questão.
Observa -se que, entre os estudantes com menor número de acertos, nessa parte do exame, a situação mais frequente foi a escolha da alternativa C ( em preto ), incorreta.
Por exemplo, entre os estudantes que acertaram três questões, 35,7 \% escolheram a alternativa C ( em preto ), 18,7 \% escolheram a alternativa E ( em vermelho ), 16,7 \% escolheram a alternativa D ( em laranja ), 15,2 \% escolheram a alternativa B ( em verde ) e 13,5 \%, a A ( em azul ).
Entre os que acertaram três respostas das questões de múltipla escolha de Formação Geral, 0,1 \% deixou a questão em branco, e 0,1 \% marcou mais de uma alternativa, invalidando a questão.
À medida que o número de acertos aumenta, indicando desempenho melhor nessa parte da prova, aumenta concomitantemente a proporção de estudantes que selecionaram a alternativa correta B, atingindo 100 \% para os estudantes com oito acertos.
Essa análise permite verificar como a questão discriminou os grupos de desempenho, justificando -se o alto índice obtido na questão.
Cumpre notar que não é possível inferir deste gráfico nem o índice de facilidade ( que seria uma média da proporção ponderada pela quantidade de alunos com cada uma das notas ), nem o índice de discriminação Ponto-Bisserial, por razão equivalente.
No caso extremo, no qual a grande concentração dos acertos dos alunos fosse abaixo de quatro, o índice de facilidade seria obrigatoriamente abaixo de 25 \% ( neste exemplo ).
Caso a concentração fosse em seis acertos ou mais, o índice seria obrigatoriamente acima de 50 \%.
Os gráficos relativos às demais questões de Formação Geral constam do Anexo I. 6.2.2 Componente de Conhecimento Específico Na Tabela 6.10, são apresentadas as Estatísticas Básicas em relação às questões objetivas do Componente de Conhecimento Específico da prova de Comunicação Social – Publicidade e Propaganda, por Grande Região.
Além disso, como será discutido ainda nesta seção, sete questões objetivas do Componente de Conhecimentos Específicos não foram usadas no cômputo das notas por terem sido descartadas pelo critério do Ponto-Bisserial.
Assim, as notas foram calculadas com base em 20 das 27 questões objetivas de Conhecimentos Específicos.
A Média do Brasil deste componente foi 55,2.
A menor Média foi observada na região Norte ( 53,3 ), e a maior, na região Sul ( 58,8 ).
O Desvio padrão de todo o Brasil foi 16,8, sendo o menor Desvio padrão encontrado na região Nordeste ( 16,3 ), e o maior, na região Sudeste ( 16,9 ).
A Mediana de todo o Brasil foi 57,1, o mesmo valor da Mediana encontrada nas regiões Nordeste, Sudeste e Centro-Oeste.
Na região Norte, a Mediana foi menor : 52,4, e, na Sul, maior, 61,9.
A nota Máxima da prova foi 100,0, obtida nas questões objetivas do Componente de Conhecimento Específico, por, pelo menos, um aluno das regiões Nordeste e Sudeste.
A nota Mínima foi zero em quase todas as regiões, exceto na região Norte, com nota Mínima 14,3.
Na Tabela 6.11, são apresentados os Índices de Facilidade e Discriminação ( Ponto-Bisserial ) das questões objetivas do Componente de Conhecimento Específico, para os estudantes de Comunicação Social – Publicidade e Propaganda.
Para facilitar a diferenciação das questões, serão usadas as mesmas cores da Tabela 6.8 para as diferentes classificações dos Índices de Facilidade e de Discriminação.
A partir do Índice de Facilidade obtido, pode -se concluir que um terço das questões objetivas da prova foram considerados, pelo menos, Difícil : das 27 questões válidas, nove foram classificadas como Difícil ou como Muito difícil.
Uma questão foi classificada como Muito Fácil, seis, como Fácil e outras 11, como Médio.
Já quanto aos Índices de Discriminação das questões objetivas do Componente de Conhecimento Específico da prova, obtém -se como resultado a seguinte classificação : seis das 27 questões foram consideradas boas, enquanto dez delas tiveram Índice de Discriminação Muito bom.
Assim, para mais da metade das questões, 16 em 27, os Índices de Discriminação foram Bom ou Muito bom.
Entre as demais, cinco delas foram classificadas como Médio, e outras seis, como Fraco, sendo 11, por conseguinte, a quantidade de questões nos dois patamares mais baixos de discriminação.
Constata -se, assim, que a prova, no que se refere ao Componente de Conhecimento Específico, possuía média capacidade de discriminar aqueles que dominam ou não o conteúdo.
O Índice de Facilidade variou de 0,11 a 0,91, e o de Discriminação, de 0,01 a 0,46.
Na Tabela 6.12, é apresentada a distribuição das questões, levando -se em conta simultaneamente a classificação dos dois índices.
Entre as questões que alcançaram os maiores índices de discriminação, dez delas foram classificadas com Índice de Discriminação Muito bom, classificadas pelo Índice de Facilidade da seguinte forma : uma como Muito Fácil, seis, como Fácil e outras três, como Médio.
As questões desse grupo classificadas como Muito Fácil ou Fácil tiveram Índice de Facilidade variando de 0,91 ( questão 22 ) a 0,63 ( questão 28 ).
Já o Índice de Facilidade das três classificadas como Médio variou entre 0,52 e 0,54.
Quanto ao Índice de Facilidade, das seis questões classificadas com Índice de Discriminação Bom, cinco foram consideradas de nível Médio e outra era Difícil.
Já das cinco questões com Índice de Discriminação Médio, três apresentavam Índice de Facilidade Médio, e duas, Difícil.
Seis questões ficaram na faixa com Índice de Discriminação Fraco, e, pelo Índice de Facilidade, classificadas como Difícil ( quatro ) e Muito Difícil ( duas ).
A questão de número 22 foi a mais fácil entre as 27 questões específicas, elevado Índice de Facilidade, 91 \% dos alunos acertaram.
Complementa a ideia o fato dessa questão apresentar poder discriminatório alto, 0,43, classificado como Muito bom, o que comprova ter sido a mais fácil para os estudantes.
Outra questão fácil foi a de número 12, com Índice de Facilidade 0,75, o que, em termos percentuais, corresponde a 75 \% de estudantes que responderam acertadamente.
O Índice de Discriminação foi 0,44, classificado, também, como Muito bom.
A questão de número 21 foi a mais difícil entre as 27 questões específicas, com baixo Índice de Facilidade, apenas 11 \% de acertos.
Essa questão apresentou poder discriminatório muito baixo, 0,02, o que comprova ter sido a mais difícil para os estudantes.
Destaca -se, também, a questão 32, com Índice de Facilidade 0,15, o que, em termos percentuais, corresponde a 15 \% de estudantes que responderam acertadamente.
Já 0,02 foi o seu Índice de Discriminação.
Tais questões foram, portanto, pelo critério Ponto-Bisserial, consideradas inadequadas.
Por isso, as questões 21 e 32, além das questões 13, 23, 25 e 34, foram eliminadas do cômputo da nota final.
A título de exemplo das análises do comportamento das questões objetivas, no Gráfico 6.5, analisa -se a questão 15 do Componente de Conhecimento Específico.
Essa questão foi considerada, pelas respostas dos estudantes avaliados na prova, como uma questão com nível de facilidade Médio, 0,54, ou seja, 54 \% dos estudantes assinalaram acertadamente a opção C, correspondente ao gabarito.
Como já comentado, seu Índice de Discriminação foi igual a 0,46, classificado como Muito bom, o maior valor de discriminação.
Neste gráfico, cada uma das cinco curvas representa o percentual de respostas em determinada alternativa da questão 15, em função do número de acertos dos estudantes nessa parte da prova, antes de possíveis eliminações de questões pelo critério do Ponto- Bisserial.
A alternativa correta C, representada no gráfico pela curva em preto, foi escolhida em maiores proporções pelos alunos com desempenho melhor nessa parte da prova.
Já as alternativas incorretas, também denominadas distratores, foram selecionadas, principalmente, por aqueles com notas mais baixas.
Observa -se que a soma não é 100 \%, por causa das questões não respondidas ou com mais de uma opção marcada.
Aqueles com nota zero, na sua quase totalidade optaram pela alternativa A ( linha azul ) e uma menor proporção deixou essa questão em branco ou marcou mais de uma alternativa, comportamento considerado inválido.
A proporção de alunos que selecionou a resposta correta C aumenta gradativamente, chegando a atingir 100 \% para 22 acertos ou mais, enquanto a proporção dos que escolheram alternativas incorretas decai, a partir de seis acertos, em função do número de acertos nessa parte da prova.
Os gráficos relativos às demais questões do Conhecimento Específico constam do Anexo I. 6.3 ANÁLISE DAS QUESTÕES DISCURSIVAS Nesta seção, apresentam -se estatísticas selecionadas e histogramas das Questões Discursivas de Formação Geral ( 6.3.1 ) e Conhecimento Específico ( 6.3.2 ).
São também apresentadas e comparadas as médias de subpopulações, caracterizadas por Grande Região. 6.3.1 Componente de Formação Geral As análises dos resultados de desempenho dos estudantes de Comunicação Social – Publicidade e Propaganda nas duas questões discursivas relativas à Formação Geral encontram -se na Tabela 6.13 e no Gráfico 6.6.
Na Tabela 6.13, observa -se que a nota Média nesse conjunto de questões foi bem abaixo da obtida nas objetivas.
Os estudantes de todo Brasil obtiveram, em Formação Geral, Média 54,9 nas questões objetivas e 38,1 nas questões discursivas.
Pode -se notar também que o Desvio padrão nesse conjunto de questões foi menor do que o obtido nas objetivas : 23,6 nas questões objetivas e 17,1 nas questões discursivas.
A maior Média foi obtida na região Sul ( 39,0 ), e a menor, na região Norte ( 36,1 ).
A Mediana de todo o Brasil, nesse componente, foi 39,0, obtida na região Sudeste.
Nas regiões Sul e Centro-Oeste, os valores da Mediana foram maiores ( 40,5 e 39,5, respectivamente ), e nas demais regiões, foram menores que o valor para o Brasil : Norte ( 36,5 ) e Nordeste ( 37,5 ).
A nota Máxima ( 93,5 ) foi obtida nas regiões Sudeste e Sul, sendo 75,5 a Máxima na região Norte ; 91,0, na região Nordeste ; e 90,0, na região Centro-Oeste.
A nota Mínima ( 0,0 ) foi a mesma em todas as regiões do Brasil.
No Gráfico 6.6, está representada a distribuição das notas nas questões discursivas do Componente de Formação Geral.
A moda dessa distribuição ocorre no intervalo ( 30 ; 40 ], seguida pelo intervalo ( 40 ; 50 ].
Destaca -se também, como um máximo local, o intervalo [ 0 ; 10 ], com distribuição de 7,6 \% do total de notas, sendo que no intervalo [ 0 ; 10 ] inclui -se, além da nota zero, a frequência de alunos que deixaram esse tipo de questão em branco.
Na sequência, os resultados verificados para cada uma das questões discursivas de Formação Geral serão apresentados, estabelecendo -se relações com os conteúdos abordados em cada uma delas.
Os comentários da Banca de docentes corretores a respeito do observado na correção das respostas dos estudantes, suas impressões e conclusões serão apresentados juntamente com a análise de cada questão.
Cumpre esclarecer que, tendo -se em vista que as questões discursivas de Formação Geral são padronizadas, ou seja, constam de todas as provas, os comentários da Banca são os mesmos para todas as carreiras acadêmicas, sendo direcionados a todos os estudantes que participaram do Enade/2018.
A seguir, serão analisados os desempenhos da Área de Comunicação Social – Publicidade e Propaganda nas duas questões discursivas de Formação Geral do Enade/2018, comparando -se os resultados obtidos e comentando -se cada questão. 6.3.1.1 Análise de Conteúdo da Questão Discursiva 1 do Componente de Formação Geral Os dados de Comunicação Social – Publicidade e Propaganda, obtidos a partir das respostas à questão 1, encontram -se na Tabela 6.14 e no Gráfico 6.7.
O desempenho médio dos estudantes na questão discursiva 1 ( média 30,2 ) foi inferior ao obtido na questão discursiva 2 ( média 34,8 ).
A maior Média para a questão 1 foi obtida na região Sul ( 30,8 ), e a menor, na região Norte ( 29,3 ).
Quanto à variabilidade das notas, o Desvio padrão de todo o Brasil foi 21,1.
O menor Desvio padrão foi obtido na região Norte ( 20,4 ), e o maior Desvio padrão foi obtido na região Sul ( 22,3 ).
As medianas de todas as regiões foram iguais a 30,0.
A nota Máxima foi 100,0 para quase todas as regiões, exceto a região Norte que teve nota Máxima 90,0.
Já a Nota Mínima foi zero para todas as regiões.
No Gráfico 6.7, mostra -se a distribuição das notas na questão discursiva 1 do Componente de Formação Geral.
Observa -se que a maior frequência ( 24,7 \% ) corresponde ao intervalo ( 30 ; 40 ].
A frequência deste intervalo modal é seguida de perto por dois outros intervalos, ( 20 ; 30 ] e ( 40 ; 50 ]. 6.3.1.2 Comentários sobre a correção de Conteúdo das respostas à Questão Discursiva 1 A Questão 1 exigia do aluno um posicionamento sobre conteúdo humanista e crítico.
O comando da questão solicitava a redação de um texto que formulasse e articulasse argumentos consistentes envolvendo temas previstos na Portaria nº 444, de 30 de maio de 2018, que dispõe sobre o componente de Formação geral do Enade 2018, tais como : solidariedade/violência, tolerância/intolerância, inclusão/exclusão, sexualidade, relações de gênero e relações \textbf{étnico-raciais} ( art. 7º, item VIII ), bem como ética, democracia e cidadania ( art. 7º, item I ).
O enunciado da questão buscava o estabelecimento da relação de três dimensões – os Direitos Humanos, a Democracia e as Fake News.
Para isso, recorreu a quatro suportes : a reprodução de uma foto e três fragmentos adaptados de textos que foram disponibilizados em páginas da Internet ( da Anistia Internacional, do Jornal El País-Brasil e do Labic-Laboratório de estudos sobre Imagem e Cibercultura ).
Os fragmentos traziam dados sobre o aumento do número de assassinatos de grupos minoritários, as origens e o trabalho da vereadora Marielle, sua luta pelos Direitos Humanos, seu assassinato, e a produção de fake news.
O comando solicitava um texto abordando dois tópicos : o tensionamento entre a defesa dos Direitos Humanos e a produção de notícias falsas, tendo como base a vereadora Marielle – a defesa dos Direitos Humanos realizada por ela e as notícias falsas divulgadas após a sua morte.
As fake news também estavam presentes no segundo tópico de análise : os prejuízos que provocam para a sociedade democrática.
Por um lado, os textos selecionados para constituir o enunciado da questão continham apoio suficiente para formulação e articulação de argumentos para discorrer sobre o tensionamento entre a atuação da Marielle em defesa dos Direitos Humanos e fake news.
Por outro, havia ênfase nos aspectos biográficos e na trajetória da vereadora Marielle, o que levou a respostas pouco focadas no comando, que discutiam a vida e a morte da vereadora, sua importância como um ícone dos Direitos Humanos, opiniões em torno do seu assassinato, até mesmo sobre a responsabilidade pelo crime, muitos textos com a presença/defesa de fake news ou produção de novas.
O foco sobre Marielle também ensejou respostas basicamente \textbf{centradas} na questão da violência.
Foram frequentes as interpretações com viés político-partidário, não previstas no padrão de resposta.
O item do comando que solicitava a discussão sobre democracia e fake news encontrou menos apoio nos textos.
Apesar de se esperar que graduandos tivessem maturidade, inclusive política, para demonstrar alguma compreensão dos aspectos principais que envolvem o convívio democrático, esse aspecto foi pouco discutido e, quando presente, quase sempre foi apenas mencionado, numa quase reprodução do comando.
A falta de entendimento do que é o conceito de Direitos Humanos também foi percebida nas respostas, seja para recomendá -lo como importante para a sociedade, seja para criticá -lo como algo que defende bandidos.
Foi marcante a ideia que vincula Direitos Humanos a direitos apenas das minorias.
Não foi incomum encontrar respostas que \textbf{personificavam} os Direitos Humanos, tratando -os como sujeito de alguma ação que deveria ser implementada na sociedade em geral ou em casos específicos.
As fake news foram citadas muitas vezes, mas pouco explicadas, conforme previsto no padrão de resposta.
O argumento utilizado geralmente vinculava as consequências de fake news para a imagem da Marielle, ou seja, centrava no fato em si, já indicado no enunciado, sem qualquer análise ou aprofundamento próprio.
É importante salientar que muitas respostas, provavelmente de grupos vinculados a cursos específicos, continham argumentações recorrentes : os Direitos Humanos como algo vinculado à Constituição e a tratados dos quais o Brasil é signatário ; as fake news como Como já apontado, houve baixa quantidade de respostas contemplando o segundo tópico do comando – prejuízos da produção de notícias falsas para a sociedade democrática ; fato motivado, possivelmente, pela falta de um elemento que ressaltasse esse ponto nos textos motivadores.
De modo geral, os estudantes não conseguiram recorrer a conhecimentos tácitos sobre democracia para construir argumentos próprios.
As respostas que abordavam o tema, geralmente estavam vinculadas à ideia do direito à verdade, do direito à integridade da pessoa ou dos direitos gerais garantidos pela Constituição.
A relação da questão com o momento político brasileiro foi indicada em um número expressivo de respostas.
Por vezes, apenas questionando a temática da questão – caso Marielle – em uma prova para avaliar o ensino superior no Brasil ou a importância dos Direitos Humanos na atual conjuntura política do país, levando à anulação da resposta.
Houve também interpretações das consequências do uso das fake news no processo eleitoral de 2018, influenciando a escolha consciente e livre das pessoas.
As respostas demonstraram atenção político social frente a esses temas, porém com pouco domínio conceitual do campo de conhecimento sobre Direitos Humanos, fake news e democracia.
Em termos da linguagem adotada pelos alunos, há dois grandes grupos de respostas.
No primeiro grupo, os que responderam de forma clara, com autonomia de pensamento e explicitação de justificativas para os aspectos solicitados no enunciado, demonstrando conhecimento sobre o que significam Direitos Humanos, fake news e democracia, além da articulação entre essas questões.
Por outro lado, um segundo grupo dos que não conseguiram ir além do que foi citado no enunciado, apenas repetindo -o, muitas vezes desviando o foco da questão para temas correlatos, como a vida de Marielle ou a violência.
Nesse segundo grupo, mais facilmente foram encontradas respostas desconexas e mal articuladas, com textos confusos que explicitam o não conhecimento sobre o tema e a falta de adequação à linguagem, necessária a um concluinte do ensino superior.
As abordagens dos Direitos Humanos mais recorrentes nas respostas corretas foram : a referência da Marielle como representante/ícone dos DH ; e os DH como direitos previstos na Constituição e em tratados internacionais dos quais o Brasil é signatário.
No entanto, também foram encontradas críticas aos Direitos Humanos, tais como : DH vistos como “ direitos de bandidos ” ; a defesa pouco embasada de frases de efeito como “ direitos humanos para humanos direitos ”.
Também chamou atenção o tratamento dos DH como um personagem, ou um órgão público, evidenciando desconhecimento sobre o tema.
Em relação às fake news, dois argumentos previstos pelo padrão de respostas foram os mais presentes : a noção de que o mundo globalizado e a facilidade de reprodução de informações pelas mídias atuais levam à criação de \textbf{mentiras} que são compartilhadas e que geram consequências ; a Marielle como alvo de fake news que tentaram desconstruir sua luta e seu legado.
Também foi frequente a ideia de que a produção de informações falsas é realizada por uma imprensa tendenciosa com intenção de manipular a informação.
No item relacionado a fake news e democracia, as respostas corretas afirmavam que as fake news inviabilizam o acesso à verdade e à garantia da integridade de um Estado Democrático de Direito.
Também foi abordada a influência das fake news nos resultados das eleições dos últimos anos, nos Estados Unidos e no Brasil, gerando consequências para a vida das pessoas.
Por fim, um outro aspecto abordado nas respostas a este tópico foi a polarização da sociedade em termos políticos, com a afirmação de que as fake news dificultam o debate, estimulam posições extremas e rivalidades prejudiciais à convivência em uma democracia plena.
A notas mais fracas ( de zero a 35 ) foram cerca de 20 \% das respostas corrigidas – as que não estavam em branco nem foram desconsideradas ou anuladas.
Foram respostas que se restringiam a cópia ou paráfrases do enunciado ; as que tinham foco em desdobramentos políticos e policiais do caso Marielle ; as que tinham argumentação confusa e/ou incompleta.
A grande maioria das notas, em torno de 65 \%, foi mediana, entre 40 a 70.
Nesse grupo ficaram as respostas que continham alguma abordagem adequada sobre Direitos Humanos, sobre as atividades da vereadora Marielle e sobre as fake news divulgadas após sua morte.
As notas mais altas deste grupo foram para respostas que incluíram a questão do tensionamento, quase sempre, em um texto coerente e consistente.
Muitas vezes, a nota mediana foi consequência do baixo índice de respostas à segunda parte do comando, que solicitava a discussão sobre os prejuízos das fake news para a democracia.
As melhores notas só foram obtidas por aqueles que incluíram a discussão sobre os prejuízos das fake news para a democracia.
Apenas 15 \% das respostas corrigidas receberam avaliação igual ou superior a 75.
A deficiência principal detectada pelos corretores foi a dificuldade de grande parte dos alunos para interpretar o enunciado da questão e para se expressar em textos consistentes.
De modo geral, os temas foram tratados de forma fragmentada e pouco analítica.
Além disso, chamou à atenção a predominância do lugar comum na construção argumentativa das respostas, a superficialidade do debate em torno da democracia e o desconhecimento de conceitos como Direitos Humanos, esses confundidos com instituições e até pessoas. 6.3.1.3 Análise de Conteúdo da Questão Discursiva 2 do Componente de Formação Geral Na Tabela 6.15, mostra -se o desempenho médio dos estudantes na questão discursiva 2.
Nessa questão – a de melhor desempenho entre as duas de Formação Geral, os alunos de todo o Brasil, obtiveram Média 34,8.
A região Sul foi aquela cuja Média, também nessa questão, foi maior ( 36,4 ), enquanto a de menor Média, como na questão discursiva 1, foi a região Norte ( 31,9 ).
Quanto à variabilidade das notas, o Desvio padrão de todo o Brasil foi 21,6, superior ao obtido na questão discursiva 1 ( 21,1 ).
O maior desvio nessa questão foi obtido na região Sul ( 23,1 ), enquanto o menor foi obtido na região Sudeste ( 21,0 ).
A Mediana de todo o Brasil foi 35,0.
Quanto às notas Máxima e Mínima, assim como na questão discursiva 1, a nota Máxima foi 100,0 para quase todas as regiões, exceto a região Norte que teve nota Máxima 90,0.
Já a Nota Mínima foi zero para todas as regiões.
No Gráfico 6.8, mostra -se a distribuição das notas na questão discursiva 2 do Componente de Formação Geral.
Observa -se que a maior frequência ( 21,4 \% ) corresponde ao intervalo ( 30 ; 40 ].
O segundo intervalo de maior frequência foi ( 40 ; 50 ].
Destaca -se, também, o intervalo ( 50 ; 60 ] com frequência de 16,1 \%.
Nota -se, ainda, que as notas ficaram ligeiramente mais dispersas em comparação às da questão discursiva de número 1, o que pode ser constatado também pela comparação do desvio padrão das notas da questão discursiva 2 ( 21,6 ) e o das notas da questão discursiva 1 ( 21,1 ). 6.3.1.4 Comentários sobre a correção de Conteúdo das respostas à Questão Discursiva 2 Os textos e o comando da Questão 2 de Formação Geral são perfeitamente adequados ao que se espera avaliar em Formação Geral, ao tratar das funções sociais dos museus, contextualizando a temática sob o ponto de vista da memória e da perspectiva de futuro.
Na questão, estavam presentes, principalmente, o conteúdo humanista e crítico, a interpretação e a produção de textos sobre temas como cultura e arte, educação, ciência e democracia, previstos na Portaria que dispõe sobre o componente de Formação Geral do Enade 2018.
Os conhecimentos exigidos para responder ao comando, explicitados no padrão de resposta, vão além do senso comum, por se exigir a compreensão de conceitos como “ função social ” ( de instituições sociais ), “ memória ” ( no sentido de memória social ), “ perspectiva de futuro ” ( no contexto \textbf{sócio-histórico-cultural} ) e “ sociedade contemporânea ” ( que deveria independer da interpretação de cada grupo social a respeito da sociedade na qual vive ).
No entanto, apesar de envolver tais conceitos, espera -se que alunos concluintes de cursos superiores de graduação sejam capazes de discutir funções de instituições sociais, especialmente porque o enunciado trazia fato recente e de grande repercussão como apoio.
O padrão de resposta elencou as funções sociais que poderiam ser abordadas nas respostas dos estudantes, como também argumentos que poderiam ser usados ao longo das mesmas.
A tarefa dos estudantes era a construção de um texto expositivo que articulasse a importância de museus, suas funções sociais e as perspectivas de memória e de futuro, o que se mostrou de dificuldade de média a difícil.
Apesar de o enunciado ser claro, muitos estudantes, provavelmente influenciados pelos textos, limitaram -se a discutir aspectos relacionados ao Museu Nacional e, em particular, ao incêndio que o \textbf{destruiu}.
Foram poucos os que interpretaram corretamente o escopo do comando, que solicitava a discussão das funções sociais de museus de uma forma geral, no sentido de argumentar a importância deste tipo de instituição.
As quatro funções sociais elencadas no padrão de resposta foram encontradas na quase totalidade das redações dos estudantes de forma articulada com a importância dos museus.
Entretanto, quase nunca os textos abordavam mais do que uma função social e vinculavam, de forma clara, a importância desse tipo de instituição para a sociedade contemporânea.
Além disso, a solicitação de contextualização sob o foco museu/memória/perspectiva de futuro parece não ter sido bem compreendida e, em consequência, esteve pouco presente nas respostas.
Também se observou confusão entre os conceitos de história e memória e uma visão estreita do significado de sociedade contemporânea, às vezes, restrita à cidade do respondente, outras, ao Rio de Janeiro, ou ainda a um tipo de sociedade idealizada.
A relevância que os textos de apoio deram ao caso particular do incêndio do Museu Nacional do Rio de Janeiro e aos prejuízos causados ao seu acervo, induziu a maioria dos estudantes a elaborar respostas tratando das funções a partir da perda e, em particular, à perda do tipo de acervo daquele museu e do tipo de pesquisa que lá se desenvolve, divergindo do comando da questão, que solicitava texto abordando funções sociais de museus em geral.
Também foram encontradas muitas respostas que se limitavam a discutir o que o estudante considerava como possíveis causas do incêndio, a falta de manutenção, questões políticas e de gestão, sem responder ao que foi solicitado.
Foi também recorrente a transcrição de frases dos textos do enunciado nas respostas, diretamente ou por paráfrases, revelando certa dificuldade de criação textual.
A análise das respostas dos estudantes, além de revelar dificuldade de abordar a temática esperada, revelou outros dois aspectos predominantes.
Um deles, como já comentado, foi a ênfase no incêndio do Museu Nacional, que orientou grande parte dos textos dos estudantes.
Outro, também evidenciando desvio do esperado, foi o foco na perda patrimonial.
Nesse grupo, houve respostas que até apresentaram uma função social, mas olhada pelo ponto de vista da perda.
Nessa linha, muitas respostas enalteciam a importância do Museu Nacional, sua relevância no cenário nacional, o impacto da perda de seu acervo, isso complementado por frágeis reflexões sobre as funções sociais.
Em consequência, a perspectiva de desdobramentos mais ricos, que tratassem de aspectos sociais relacionados a diferentes tipos de museu foi prejudicada.
Muitas respostas trataram da preservação da memória e/ou guarda do acervo ; identidade, tanto individual quanto coletiva ; conhecimento sobre antepassados ; evolução da sociedade e dos seus costumes ; valorização e disseminação do conhecimento por meio da liberação do acesso aos acervos ; bem como atividades pedagógicas e exposições de arte.
O aspecto produção de pesquisa e conhecimento, embora frequente, foi mencionado de forma muito simples, sem qualquer aprofundamento, para além do que já era mencionado no enunciado.
Frequentemente apresentava -se o museu como produtor de pesquisas para o campo da saúde, o que evidencia desconhecimento das reais possibilidades de pesquisa, mesmo no Museu Nacional.
Quanto a perspectiva de futuro, quando presente no texto, era mencionada de forma breve, com afirmações genéricas sobre a possibilidade de melhoria do futuro pelo simples conhecimento do passado, numa percepção de que o passado determina o futuro, de forma linear.
Foram raras as respostas com alguma argumentação que associasse pesquisa, memória, identidade e reflexão social como base para a construção de alternativas de futuro.
A relação entre turismo e museu apareceu considerando a visita a museus como forma de lazer e, algumas vezes, como geradora de recursos.
O desempenho pode ser considerado fraco.
Muitas respostas deixavam claro que o estudante nunca tinha visitado um museu e, de uma forma geral, observou -se desconhecimento sobre suas funções sociais.
Não se esperava que história, memória social e museologia fossem assuntos dominados pelos estudantes das diversas áreas que participaram do Enade 2018, mas esperava -se que essa temática pudesse ser abordada do ponto de vista das experiências pessoais ou de conhecimento social.
Como uma vivência própria com os museus também se revelou frágil ou inexistente, o conhecimento sobre o assunto ficou bem restrito aos textos motivadores e ao senso comum, construído, basicamente, pelo noticiário em torno do incêndio.
O universo dos argumentos dos estudantes se mostrou muito limitado.
Cada dimensão foi praticamente representada por uma ideia simplista.
De forma recorrente : a preservação da memória para construção da identidade de um povo ; a compreensão do passado por meio das pesquisas como oportunidade para melhorar o futuro ; o museu como um espaço em que há um \textbf{aprendizado} estimulado pelo lazer.
Houve, também, uma grande concentração de respostas associando o museu a um local de guarda de acervo e de história.
As respostas à Questão 2 evidenciaram dificuldades conceituais.
Além dos conceitos de história, memória e museu se misturarem, parecendo não se distinguirem, ao mencionarem a pesquisa, também se observou falta de distinção entre pesquisas colegiais, com objetivos de \textbf{aprendizado}, e as realizadas por pesquisadores, com objetivos de avançar cientificamente.
A falta de compreensão do que se pedia no comando da questão e a dificuldade de expressão linguística foram dois aspectos marcantes na avaliação dos corretores.
Foi frequente os textos apresentarem ideias soltas ou de forma confusa, sem constituir unidade.
O vocabulário usado se revelou bem restrito, e fortemente apoiado no coloquial, na linguagem oral.
Predominou o senso comum em lugar de uma reflexão mais formal envolvendo os conceitos envolvidos. 6.3.1.5 Análise de Língua Portuguesa das Questões Discursivas do Componente de Formação Geral Os dados de Comunicação Social – Publicidade e Propaganda, obtidos a partir das respostas às questões discursivas do Componente de Formação Geral, no que tange à Língua Portuguesa, encontram -se na Tabela 6.16 e no Gráfico 6.9.
Nesse aspecto, os alunos de todo o Brasil obtiveram Média 60,7.
A maior Média em relação à Língua Portuguesa foi obtida na região Sudeste ( 61,7 ), e a menor, na região Nordeste ( 57,2 ).
Quanto à variabilidade das notas, o Desvio padrão de todo o Brasil foi 20,4.
O menor Desvio padrão foi obtido na região Sudeste ( 19,4 ), e o maior Desvio padrão foi obtido na região Centro-Oeste ( 23,3 ).
A Mediana das notas de Língua Portuguesa foi 67,5 para o Brasil, como um todo, a mesma obtida nas regiões Sudeste e Sul.
Nas demais regiões, a Mediana foi 65,0.
A nota Máxima para todo o Brasil foi 92,5, com, pelo menos, um aluno tirando essa nota nas regiões Sudeste e Centro-Oeste.
Na região Norte, a nota Máxima foi 87,5, e nas regiões Nordeste e Sul, foi 90,0.
Já a nota Mínima foi zero em todas as regiões do país.
No Gráfico 6.9, mostra -se a distribuição das notas de Língua Portuguesa do Componente de Formação Geral.
Observa -se que a maior frequência ( 32,1 \% ) corresponde à dos alunos que obtiveram nota no intervalo ( 60 ; 70 ].
Destacam -se, também, os alunos que que obtiveram nota no intervalo ( 70 ; 80 ), correspondendo a 28,1 \%.
Os alunos que deixaram ambas as questões em branco representam um pouco mais do que 5 \% do total, caracterizando -se como máximo local. 6.3.1.6 Comentários sobre a correção das respostas de Formação Geral com respeito à Língua Portuguesa As duas questões discursivas de Formação Geral ensejaram duas configurações textuais distintas, tendo em vista os seus encaminhamentos : a Questão 1 teve um caráter mais opinativo e a Questão 2, mais expositivo, conforme se explicita a seguir.
Ambas as questões permitiram a \textbf{observação} do desempenho linguístico dos participantes em situação formal, competência exigida no processo de formação universitária.
As habilidades de produção textual envolvidas têm, igualmente, como condição básica, o domínio do padrão formal culto da Língua Portuguesa, no que diz respeito aos aspectos relativos às questões de caráter ortográfico, textual, morfossintático e vocabular.
O encaminhamento das questões evidencia a expectativa de que o participante utilizasse seus conhecimentos sobre os dois temas e estruturasse seus textos de acordo com as características do registro formal adequado à situação comunicativa.
Tal configuração determina exigências quanto : à adequação da seleção vocabular, ao desenvolvimento do conteúdo, à estruturação sintática dos períodos, à organização lógica das ideias, à utilização de procedimentos de encadeamento textual e de referenciação, à obediência às exigências morfossintáticas próprias da modalidade escrita da norma-padrão, ao respeito às regras ortográficas e às regras de acentuação gráfica.
Esse encaminhamento exigiu do graduando o domínio do modo de organização textual de base dissertativo-argumentativa, essencial ao processo de formação universitária.
As habilidades de produção textual envolvidas têm como condição básica o domínio do padrão formal culto da Língua Portuguesa, no que diz respeito aos aspectos relativos às questões de caráter ortográfico, textual, morfossintático e vocabular.
O enunciado da Questão 1 tomava como base três fragmentos de textos jornalísticos publicados em sites da internet, que abordavam o aumento do assassinato de minorias e de defensores de direitos humanos no Brasil, entre eles a vereadora Marielle Franco.
O comando solicitava que o participante desenvolvesse um texto dissertativo e argumentativo com as temáticas Fake News, Democracia e Direitos Humanos.
Embora não tenha havido referência explícita à tipologia textual esperada para o desenvolvimento da resposta, essa proposta encaminhou o participante para a elaboração de um texto opinativo, com o objetivo de discutir o problema e manifestar um posicionamento crítico.
O enunciado da Questão 2 abordava o incêndio do Museu Nacional do Rio de Janeiro, ocorrido em 2018, e seus impactos na ciência brasileira e internacional, tomando como base dois fragmentos de textos jornalísticos publicados em sites da internet.
O comando solicitava que o participante desenvolvesse um texto abordando três aspectos da função social dos museus na sociedade contemporânea, encaminhando -o para a elaboração de um texto expositivo que confirmasse a posição defendida nos fragmentos motivadores a respeito da importância dessas instituições.
Para atender a essas exigências, o padrão de resposta utilizado na avaliação das questões discursivas de Formação Geral considerou aspectos relevantes ao bom desempenho linguístico como competências distintas, de modo a permitir um mapeamento detalhado do domínio dos recursos disponíveis na Língua Portuguesa para a comunicação escrita formal : ortográficos, textuais, morfossintáticos e vocabulares ( ver detalhes no Anexo VIII – Padrão de Respostas ).
Aspectos Ortográficos Para avaliar a competência relativa ao domínio das convenções ortográficas da norma-padrão da Língua Portuguesa, observou -se o respeito às regras de acentuação gráfica e da grafia padrão das palavras ( com ausência de abreviaturas próprias da linguagem da internet ), de acordo com as convenções estabelecidas pela legislação em vigor e consubstanciadas no Vocabulário Ortográfico da Língua Portuguesa, editado pela Academia Brasileira de Letras.
Aspectos Textuais Para avaliar a competência relativa ao domínio dos procedimentos de estruturação textual, procurou -se observar a utilização de mecanismos linguístico-discursivos que estabeleçam o encadeamento lógico entre as partes do texto, de forma a garantir a progressão e a coerência textuais.
Esta competência destaca -se por sua relevância, já que são avaliadas a clareza e a unidade textual.
Aspectos Morfossintáticos Para avaliar o domínio dos diferentes aspectos morfossintáticos próprios da modalidade escrita formal da norma-padrão da Língua Portuguesa, observou -se se o texto produzido atendeu às seguintes exigências : a concordância nominal, a concordância verbal, a regência nominal, a regência verbal, a flexão nominal, a flexão verbal, a correlação entre os tempos verbais, a colocação pronominal e a utilização de sinais de pontuação que contribuam para a organização lógica da frase e do texto.
Aspectos Vocabulares Para avaliar a adequação da seleção vocabular à modalidade escrita formal da Língua Portuguesa, exigida pela situação comunicativa, observou -se o atendimento às seguintes exigências : precisão na seleção/utilização do vocabulário relacionado à temática solicitada pela questão ; ausência de marcas da oralidade, como termos de sentido muito genérico ( “ coisa ”, “ negócio ”, “ você ” ) e termos de registros mais informais ( como gírias, jargões, frases feitas, ditados populares, termos regionais ).
Assim, espera -se que o participante respeite a adequação vocabular não usando gírias ou expressões coloquiais, evite repetição desnecessária de palavras e utilize um vocabulário mais formal, como solicitado por um texto dissertativo.
Para efeito de pontuação, a avaliação do desempenho linguístico considerou três grandes grupos de competências, agrupando aspectos explicitados, com os seguintes \textbf{pesos} relativos : aspectos ortográficos ( 20 \% ) ; aspectos textuais ( 40 \% ), aspectos morfossintáticos e vocabulares ( 40 \% ).
Os critérios de pontuação utilizados no padrão de resposta procuraram atender à diversidade de desempenho dos graduandos.
Nos textos analisados, observaram -se, assim, vários níveis de atendimento às exigências do padrão formal próprio da modalidade escrita da Língua Portuguesa, configurando um"continuum".
Por um lado, textos que traduzem com exatidão a expectativa de uso formal da língua, em todos os níveis analisados, marcados por clareza e unidade ; e, no outro extremo, textos quase incompreensíveis devido ao grande número de desvios de todos os tipos, principalmente os de caráter textual.
As duas questões discursivas de Formação Geral analisadas revelaram desempenhos distintos dos participantes.
A Questão 1 propiciou a oportunidade de maior desenvolvimento do tema solicitado, gerando consequências na estruturação textual, já que os textos foram mais longos, mais elaborados e fluentes.
Em virtude da polêmica sobre a existência das fake news e sobre a visão pública da personagem da vereadora Marielle Franco, observaram -se respostas que polemizaram o tema, favorecendo a construção argumentativa do texto.
A Questão 2 teve um comportamento distinto, devido ao caráter mais expositivo do enunciado e ao reduzido conhecimento dos participantes sobre o tema.
Assim, os textos foram mais curtos, frequentemente com repetição de ideias dos textos motivadores.
Quanto aos aspectos linguísticos analisados durante esta avaliação, observaram -se os seguintes resultados : Aspectos ortográficos : O desempenho dos participantes revelou uma diferença muito grande nos dois aspectos analisados nesta competência : baixo índice de desvios da grafia padrão e grande índice de desvios de acentuação.
Em vários casos, ocorre ausência completa de acentuação gráfica em todas as palavras do texto.
Os resultados revelaram que a tendência dominante entre os universitários brasileiros é a eliminação da acentuação gráfica, talvez motivada pelos hábitos relacionados às redes sociais e pela ausência de esclarecimento dos meios de comunicação, das autoridades e das escolas sobre as decisões do Acordo Ortográfico da Língua Portuguesa de 1990.
Os casos mais sistemáticos de eliminação do acento indicador da sílaba tônica foram : - palavras proparoxítonas ( por exemplo “ generos ”, “ arbitro ”, “ politico ”, “ publicas ”, “ numero ” ) ; - palavras paroxítonas terminadas em ditongo crescente ( por exemplo “ area ”, “ noticiarios ”, “ individuo ”, “ historia ”, “ varios ”, “ noticia ”, “ negligencia ”, “ ocorrencia ” ) ; - palavras oxítonas ( por exemplo “ ninguem ”, “ esta ”, “ ate ”, “ tambem ” ).
Por outro lado, destaca -se o uso indevido do acento gráfico em palavras como “ melâncolia ” e “ intervênção ”, por exemplo.
Quanto ao domínio das convenções relativas à grafia das palavras, observaram -se desvios \textbf{esporádicos} como : hipercorreção pela escolha de “ e ” no lugar de “ i ”, por influência de hábitos da oralidade ( “ descriminação ” no lugar de “ discriminação ”, “ entervenção ” no lugar de “ intervenção ” ) ; desvios de grafia relacionados à variação diastrática, como por exemplo “ perca ”, “ subjulgadas ”, “ precoseito ”, “ soubre ”, “ vecendo ”, “ indesencia ”, “ apolojia ”, “ fulga ”, “ dereitos ” ;...
Observaram -se, também, muitos casos de inadequação no uso da maiúscula : uso indevido para destacar determinadas palavras-chave do texto, como “ Brasileiros ”, “ Homens ” ; ausência de maiúscula para grafar nomes próprios, como"marielle".
Expressões como"Direitos Humanos"foram grafadas diversificadamente, com maiúsculas e com minúsculas.
Destaque -se, também, a maioria dos participantes que grafam os textos inteiramente em caixa alta, o que impede a identificação de desvios relacionados ao uso de maiúsculas.
Vale observar, também, que, ao contrário do que se esperava, não apareceram abreviaturas próprias do “ internetês ”, relacionadas ao uso de redes sociais e e-mails, como por exemplo"vc".
Aspectos textuais : Para alguns dos graduandos essa competência se revelou como a mais problemática, tendo em vista os inúmeros problemas observados, desvios acumulados durante toda a formação escolar.
São eles : sequência justaposta de ideias sem encaixamentos sintáticos, provocando uma redução de estruturas subordinadas, ao lado do aumento na frequência de estruturas coordenadas e absolutas ; frases fragmentadas que comprometem a estrutura lógico-gramatical ; frases formadas apenas por oração subordinada, sem oração principal ; redução no uso de conectores para expressar relações lógicas essenciais à construção do texto como consequência da mudança de estruturação frasal ; emprego inadequado do pronome relativo ( com omissão da preposição ou a utilização do pronome relativo inadequado, como “ onde ” ), refletindo hábitos da oralidade ; ausência de recursos de referenciação, como substituição de termos por sinônimos, hiperônimos, nominalizações, expressões metafóricas.
Esses problemas revelam dificuldades em relação à estrutura formal do texto produzido, o que é preocupante ao se levar em conta que são graduandos em fase final de formação.
Em alguns textos, falta um mínimo de textualidade e de domínio do registro padrão da língua.
Na verdade, observaram -se relações linguísticas quase agramaticais, como as estabelecidas pela sequência de gerúndios sem o apoio de um ponto de partida para a organização das informações gramaticais e semânticas.
Quanto à utilização dos sinais de pontuação, observou -se uma ausência dos mesmos na maioria dos textos analisados.
Foi muito frequente a ocorrência de parágrafos sem marca interna de pontuação para separar os períodos.
São os seguintes os tipos de problemas mais encontrados :  vírgula : utilização de vírgula para separar o sujeito e o predicado ; uso de vírgula no lugar do ponto para separar ideias que constituem períodos distintos ; ausência de vírgula para separar oração adjetiva explicativa ;  ponto e vírgula : utilização do ponto e vírgula no lugar de vírgula ;  ponto final : ausência de ponto final para separar períodos.
Aspectos morfossintáticos e vocabulares : Os resultados são transparentes em relação aos aspectos mais problemáticos no desempenho dos participantes nestas duas competências.
Em relação aos aspectos morfossintáticos, seguem algumas \textbf{observações} sobre os desvios mais frequentes.
Quanto à regência, o desvio mais frequente foi a falta do sinal indicativo da crase – isso revela que o usuário não tem consciência de que, sob a forma do termo “ a ”, existe a presença de uma contração entre a preposição “ a ” ( exigida pela regência do termo anterior ) e o artigo definido “ a ”.
Outro problema relacionado à regência verbal e à nominal, encontrado frequentemente nas respostas, foi a ausência de preposição antes de pronome relativo, processo generalizado lógico-gramatical ; frases formadas apenas por oração subordinada, sem oração principal ; redução no uso de conectores para expressar relações lógicas essenciais à construção do texto como consequência da mudança de estruturação frasal ; emprego inadequado do pronome relativo ( com omissão da preposição ou a utilização do pronome relativo inadequado, como “ onde ” ), refletindo hábitos da oralidade ; ausência de recursos de referenciação, como substituição de termos por sinônimos, hiperônimos, nominalizações, expressões metafóricas.
Esses problemas revelam dificuldades em relação à estrutura formal do texto produzido, o que é preocupante ao se levar em conta que são graduandos em fase final de formação.
Em alguns textos, falta um mínimo de textualidade e de domínio do registro padrão da língua.
Na verdade, observaram -se relações linguísticas quase agramaticais, como as estabelecidas pela sequência de gerúndios sem o apoio de um ponto de partida para a organização das informações gramaticais e semânticas.
Quanto à utilização dos sinais de pontuação, observou -se uma ausência dos mesmos na maioria dos textos analisados.
Foi muito frequente a ocorrência de parágrafos sem marca interna de pontuação para separar os períodos.
São os seguintes os tipos de problemas mais encontrados : • vírgula : utilização de vírgula para separar o sujeito e o predicado ; uso de vírgula no lugar do ponto para separar ideias que constituem períodos distintos ; ausência de vírgula para separar oração adjetiva explicativa ; • ponto e vírgula : utilização do ponto e vírgula no lugar de vírgula ; • ponto final : ausência de ponto final para separar períodos.
Aspectos morfossintáticos e vocabulares : Os resultados são transparentes em relação aos aspectos mais problemáticos no desempenho dos participantes nestas duas competências.
Em relação aos aspectos morfossintáticos, seguem algumas \textbf{observações} sobre os desvios mais frequentes.
Quanto à regência, o desvio mais frequente foi a falta do sinal indicativo da crase – isso revela que o usuário não tem consciência de que, sob a forma do termo “ a ”, existe a presença de uma contração entre a preposição “ a ” ( exigida pela regência do termo anterior ) e o artigo definido “ a ”.
Outro problema relacionado à regência verbal e à nominal, encontrado frequentemente nas respostas, foi a ausência de preposição antes de pronome relativo, processo generalizado na modalidade oral da língua, em situações de registro informal.
Apesar da possibilidade de que essa alteração de regência se generalize no padrão escrito da Língua Portuguesa, como já está ocorrendo até em textos jornalísticos, o não emprego da preposição foi considerado como desvio neste processo de avaliação.
Outro desvio de regência significativo foi a utilização inadequada de uma preposição ou sua ausência após o verbo ou o nome ( substantivo ou adjetivo ).
A concordância verbal e a concordância nominal apresentaram alguns desvios muito frequentes.
Quanto à concordância de número, observou -se, como apontam as pesquisas nessa área, ausência de marca ( com sujeito anteposto ou posposto ) ou uso indevido ( uso inadequado da marca de plural comandado pelo núcleo plural da locução adjetiva, apesar de o substantivo que funciona como núcleo do sintagma nominal estar no singular ).
Uma ocorrência que se destacou foi a ausência de acento circunflexo na forma plural do presente do indicativo do verbo “ ter ”, que foi considerada como um desvio na concordância verbal e não na acentuação gráfica.
Quanto à concordância de gênero, vários casos foram observados, normalmente no âmbito de sintagmas nominais longos, em que o adjetivo está afastado do substantivo.
Quanto à questão da colocação pronominal, foram poucos os desvios observados.
Concluiu -se que, no registro escrito formal, a maioria dos participantes já incorporou regras como a não introdução da frase por um pronome oblíquo e a próclise na presença de um termo atrator.
Várias marcas de oralidade foram identificadas, embora não com alta frequência : uso do pronome relativo “ onde ” como relativo universal, falta de artigo definido antes de substantivo isolado ou antes de termos em enumerações.
Quanto aos aspectos vocabulares, alguns tipos de inadequação foram observados : expressões da oralidade ; seleção vocabular incompatível com o contexto, gerando falta de inteligibilidade ; falta de domínio de vocabulário mais abstrato e de maior complexidade, essencial ao desenvolvimento do texto de base dissertativa.
O principal aspecto observado foi a excessiva repetição de certas palavras, como o termo"pessoa", por exemplo, revelando limitação de repertório vocabular.
Em função do tema solicitado na Questão 1, diferentes grafias da expressão fake news não foram consideradas.
Para sintetizar, os problemas que mais se destacaram nos textos analisados foram : a ) truncamentos sintáticos, caracterizados por estruturas frasais incompreensíveis devido à ruptura da complexidade sintática própria da modalidade escrita ; na modalidade oral da língua, em situações de registro informal.
Apesar da possibilidade de que essa alteração de regência se generalize no padrão escrito da Língua Portuguesa, como já está ocorrendo até em textos jornalísticos, o não emprego da preposição foi considerado como desvio neste processo de avaliação.
Outro desvio de regência significativo foi a utilização inadequada de uma preposição ou sua ausência após o verbo ou o nome ( substantivo ou adjetivo ).
A concordância verbal e a concordância nominal apresentaram alguns desvios muito frequentes.
Quanto à concordância de número, observou -se, como apontam as pesquisas nessa área, ausência de marca ( com sujeito anteposto ou posposto ) ou uso indevido ( uso inadequado da marca de plural comandado pelo núcleo plural da locução adjetiva, apesar de o substantivo que funciona como núcleo do sintagma nominal estar no singular ).
Uma ocorrência que se destacou foi a ausência de acento circunflexo na forma plural do presente do indicativo do verbo “ ter ”, que foi considerada como um desvio na concordância verbal e não na acentuação gráfica.
Quanto à concordância de gênero, vários casos foram observados, normalmente no âmbito de sintagmas nominais longos, em que o adjetivo está afastado do substantivo.
Quanto à questão da colocação pronominal, foram poucos os desvios observados.
Concluiu -se que, no registro escrito formal, a maioria dos participantes já incorporou regras como a não introdução da frase por um pronome oblíquo e a próclise na presença de um termo atrator.
Várias marcas de oralidade foram identificadas, embora não com alta frequência : uso do pronome relativo “ onde ” como relativo universal, falta de artigo definido antes de substantivo isolado ou antes de termos em enumerações.
Quanto aos aspectos vocabulares, alguns tipos de inadequação foram observados : expressões da oralidade ; seleção vocabular incompatível com o contexto, gerando falta de inteligibilidade ; falta de domínio de vocabulário mais abstrato e de maior complexidade, essencial ao desenvolvimento do texto de base dissertativa.
O principal aspecto observado foi a excessiva repetição de certas palavras, como o termo"pessoa", por exemplo, revelando limitação de repertório vocabular.
Em função do tema solicitado na Questão 1, diferentes grafias da expressão fake news não foram consideradas.
Para sintetizar, os problemas que mais se destacaram nos textos analisados foram : b ) redução de estruturas subordinadas, compensada pelo aumento na frequência de estruturas coordenadas e absolutas, por um lado, ou pela elaboração de estruturas textuais prejudicadas pelo excesso de ideias sem a devida conexão subordinativa ; c ) redução no uso de conectores para expressar relações lógicas essenciais à construção do texto, substituídos pela exigência de inferência por parte do interlocutor para suprir a sua ausência ; d ) redução no uso do subjuntivo, ao lado da ampliação do uso do indicativo, combinado a estruturas frasais coordenadas ou absolutas ; d ) empobrecimento do processo de referenciação, com a repetição de palavras, sem processos mais sofisticados de substituição ; e ) simplificação da marcação da categoria tempo na morfologia verbal ; f ) desrespeito às regras de concordância verbal e nominal próprias do padrão formal ; g ) ausência do acento grave indicador da crase, revelando falta de conhecimento dos critérios de regência verbal e nominal ; h ) falta de domínio de vocabulário mais abstrato e de maior complexidade, essencial ao desenvolvimento do processo dissertativo ; i ) redução do emprego da acentuação gráfica.
Como resultado da correção, merece atenção a constatação de que a modalidade escrita tem apresentado a tendência a uma evidente simplificação, aproximando -se das características da modalidade oral da Língua Portuguesa.
No caso dos textos de base dissertativa ( expositiva ) e de base argumentativa, inscritos em um registro formal, a distância entre as duas modalidades é ainda maior, o que provoca desvios recorrentes em todos os quatro aspectos analisados : ortográficos, textuais, morfossintáticos e vocabulares.
Destaca -se, nessa avaliação, o comprometimento textual, por meio de estruturas fragmentadas e/ou truncadas, rompendo a complexidade sintática esperada no padrão formal. 6.3.2 Componente de Conhecimento Específico Na parte da prova relativa às questões discursivas do Componente de Conhecimento Específico ( Tabela 6.17 ), observa -se que a Média foi mais baixa do que para as questões discursivas do Componente de Formação Geral.
Enquanto no Componente de Formação Geral, a Média para estudantes de Comunicação Social – Publicidade e Propaganda de todo o Brasil foi 38,1, na parte de Conhecimento Específico, a Média foi 24,3.
A maior Média deste componente foi obtida pelos estudantes da região Sul ( 26,3 ), e a menor, pelos da região Norte ( 20,9 ).
Quanto à variabilidade das notas, o Desvio padrão de todo o Brasil foi 14,9.
O maior Desvio padrão foi encontrado na região Sul ( 15,6 ), e o menor, na região Norte ( 13,4 ).
A maior nota Máxima, 80,0, foi obtida na região Sudeste.
Nas demais regiões, as notas Máximas foram : 60,0 na região Norte, 73,3 na região Nordeste, 76,7 na região Sul e 66,7 na região Centro-Oeste.
Já a nota Mínima ( 0,0 ) foi obtida por, pelo menos, um aluno em todas as regiões do Brasil.
A Mediana do Brasil, como um todo, foi 23,3.
A maior Mediana foi obtida na região Sul ( 26,7 ), e a menor, nas regiões Norte e Nordeste ( 20,0 ).
No Gráfico 6.10, representa -se a distribuição das notas nas questões discursivas no Componente de Conhecimento Específico.
A moda dessa distribuição ocorre no intervalo [ 20 ; 30 ], com mais do que 24 \% do total de participantes.
Esse intervalo é seguido de perto pelo intervalo [ 0 ; 10 ], no qual estão computadas as ocorrências de respostas em branco para as três questões discursivas de Conhecimentos Específicos, e também pelo intervalo [ 10 ; 20 ], apresentando, ambos, frequência igual 21,9 \%.
A análise de cada uma dessas questões será feita a seguir. 6.3.2.1 Análise da Questão Discursiva 3 do Componente de Conhecimento Específico Na questão 3, cujos resultados aferidos encontram -se descritos na Tabela 6.18, a Média dos estudantes de todo o Brasil foi 23,7.
O desempenho dos estudantes nessa questão foi inferior ao da questão 4 e superior ao da questão 5.
A menor Média nessa questão foi obtida pelos alunos da região Norte ( 19,7 ), enquanto a maior Média foi obtida pelos alunos da região Sul ( 27,2 ).
Quanto à variabilidade das notas, o Desvio padrão de todo o Brasil foi 19,2.
O maior Desvio padrão foi obtido na região Centro-Oeste ( 20,1 ), enquanto o menor foi obtido na região Norte ( 18,8 ).
A nota Máxima, 100,0 pontos, foi a mesma obtida nas regiões Sudeste e Centro-Oeste.
Nas regiões Norte e Nordeste, a Máxima foi 80,0, e na região Sul, foi 90,0.
A Mediana do Brasil, como um todo, foi 20,0, a mesma para quase todas as regiões, com exceção da região Sul, na qual a Mediana foi 30,0.
A nota Mínima ( 0,0 ) foi a mesma em todas as regiões do Brasil.
No Gráfico 6.11, mostra -se a distribuição das notas na questão discursiva 3, do Componente de Conhecimento Específico da área de Comunicação Social – Publicidade e Propaganda.
Essa é uma distribuição bimodal, com a moda principal no intervalo de notas ( 20 ; 30 ], com frequência de 18,4 \%.
A moda secundária situa -se na nota zero, com frequência de 18,2 \%.
Destaca -se o intervalo ( 10 ; 20 ] com frequência muito próxima à modal, 17,2 \%. 6.3.2.2 - Comentários sobre as respostas à Questão Discursiva 3 A Questão 3 demandava conhecimentos sobre técnicas de convencimento do público e atração para a marca e incluía a dimensão ética do trabalho da publicidade, na medida em que ressaltava a sua função social de dar visibilidade para grupos minoritários, valorizar a diversidade, quebrar estereótipos e incentivar a adesão de causas sociais.
O enunciado trazia um texto de opinião, publicado em uma página sobre comunicação, marketing e mídia do país ( www.meioemensagem.com.br ), abordando um comercial da Natura para a TV aberta, que envolvia uma simpatia ( “ amarrar a pessoa amada ” ) e, ao final, buscava surpreender os espectadores.
Já o comando trazia o uso do humor para tratar de temáticas sensíveis na sociedade brasileira.
A questão solicitava a análise de desenvolvimento de soluções criativas, em relação a estéticas e linguagens, além do exercício ético da profissão.
Desse modo, abordava temas presentes nas Diretrizes Curriculares da Área.
No entanto, é importante ressaltar que as teorias do riso não são itens obrigatórios das Diretrizes Curriculares de Publicidade e Propaganda, o que pode ter levado os estudantes a darem respostas mais intuitivas e especulativas.
Além disso, o exemplo usado não é diretamente uma situação de humor, fato que levou a interpretações diferentes da esperada pelo padrão de resposta.
Também houve dificuldade na realização de análises das questões éticas, a partir de um ponto de vista profissional e não pessoal.
Apontar os riscos e as oportunidades de uma marca investir na criação publicitária, a princípio, levaria a respostas diretas, objetivas e organizadas.
No entanto, fazer tal análise, no caso de utilização de estratégias de humor, como demandava a questão, não era simples.
Para atingir essa clareza, teria sido útil apresentar a definição do conceito de humor a ser trabalhada, como também de um caos/exemplo que recorresse efetivamente ao humor.
A questão pode ser considerada de nível médio em comparação com as demais questões de formação específica.
De um modo geral, as respostas não contemplaram os pontos principais demandados pela questão.
Uma grande parte dos estudantes focou em um discurso opinativo sobre questões éticas na abordagem de temas que poderiam ser problemáticos, se compreendidos de forma pejorativa.
Foram dadas muitas respostas baseadas em opiniões baseadas em discursos religiosos, identitários ou moralistas.
Como não houve uma definição do conceito de humor a ser trabalhado, grande parte das respostas deixaram de fora as possibilidades do uso do humor em favor da marca.
A abordagem mais usada para responder à questão foi emitir opiniões de que o humor poderia ser considerado pejorativo quando usado para tratar causas relacionadas a questões sociais, de gênero ou religiosas.
O conceito de humor foi usado, na maioria dos casos, como algo que pudesse diminuir o outro e não como uma estratégia que pudesse dar leveza e atrair a atenção do público para a marca.
É possível que o caso exemplificado no enunciado tenha direcionado a atenção dos alunos para esse tipo de abordagem.
De todo modo, ficaram de fora conhecimentos sobre técnicas de convencimento do público e atração para a marca, pontos fundamentais exigidos pela questão.
A linguagem adotada pelos estudantes foi menos crítica e analítica, como seria esperado de um texto de análise acadêmica, e mais coloquial, o que pode ser reflexo de abordagens de análises mais pessoais.
As respostas revelaram, ainda, um péssimo uso das normas da língua portuguesa escrita, tanto do ponto de vista da gramática quanto da ortografia.
Por conta da escolha do ponto de vista, também houve diversidade nas respostas.
Por vezes, os estudantes responderam a partir do ponto de vista dos consumidores ( em sua maioria ), denunciando uma realidade da qual discordam.
Uma minoria se colocou a partir do ponto de vista do profissional ou ainda do universitário que reflete criticamente sobre a questão.
Ficou evidente que faltou articulação de ideias e entendimento da importância da reflexão sobre os modos de atuar eficazmente e, ao mesmo tempo, eticamente no mercado.
A média das notas ficou abaixo de 25 pontos, um quarto da pontuação máxima.
Houve grande quantidade ( 65 \% dos estudantes ) de notas baixas, entre zero e 30.
Além disso, quase 7 \% dos participantes deixaram a resposta em branco, também recebendo nota zero.
O percentual de estudantes cujas respostas atingiram um patamar mediano de nota, entre 40 a 70, foi de 28 \%.
Uma minoria, menos de 1 \% do total de participantes, conseguiu notas boas, acima de 75.
As respostas encontradas revelam que há algumas deficiências na formação da área, como a importância do desenvolvimento da análise crítica e articulada sobre questões com as quais os alunos vão se deparar no mercado de trabalho.
Também ficou evidente que grande parte dos formandos não sabe diferenciar uma opinião pessoal de uma análise crítica, que aborde a complexidade do tema. 6.3.2.3 - Análise da Questão Discursiva 4 do Componente de Conhecimento Específico Na Tabela 6.19, constam as informações relativas à questão 4 do conjunto de questões do Componente de Conhecimento Específico.
O desempenho dos estudantes de todo o Brasil nesta questão foi o mais alto entre as três questões discursivas desse de gênero ou religiosas.
O conceito de humor foi usado, na maioria dos casos, como algo que pudesse diminuir o outro e não como uma estratégia que pudesse dar leveza e atrair a atenção do público para a marca.
É possível que o caso exemplificado no enunciado tenha direcionado a atenção dos alunos para esse tipo de abordagem.
De todo modo, ficaram de fora conhecimentos sobre técnicas de convencimento do público e atração para a marca, pontos fundamentais exigidos pela questão.
A linguagem adotada pelos estudantes foi menos crítica e analítica, como seria esperado de um texto de análise acadêmica, e mais coloquial, o que pode ser reflexo de abordagens de análises mais pessoais.
As respostas revelaram, ainda, um péssimo uso das normas da língua portuguesa escrita, tanto do ponto de vista da gramática quanto da ortografia.
Por conta da escolha do ponto de vista, também houve diversidade nas respostas.
Por vezes, os estudantes responderam a partir do ponto de vista dos consumidores ( em sua maioria ), denunciando uma realidade da qual discordam.
Uma minoria se colocou a partir do ponto de vista do profissional ou ainda do universitário que reflete criticamente sobre a questão.
Ficou evidente que faltou articulação de ideias e entendimento da importância da reflexão sobre os modos de atuar eficazmente e, ao mesmo tempo, eticamente no mercado.
A média das notas ficou abaixo de 25 pontos, um quarto da pontuação máxima.
Houve grande quantidade ( 65 \% dos estudantes ) de notas baixas, entre zero e 30.
Além disso, quase 7 \% dos participantes deixaram a resposta em branco, também recebendo nota zero.
O percentual de estudantes cujas respostas atingiram um patamar mediano de nota, entre 40 a 70, foi de 28 \%.
Uma minoria, menos de 1 \% do total de participantes, conseguiu notas boas, acima de 75.
As respostas encontradas revelam que há algumas deficiências na formação da área, como a importância do desenvolvimento da análise crítica e articulada sobre questões com as quais os alunos vão se deparar no mercado de trabalho.
Também ficou evidente que grande parte dos formandos não sabe diferenciar uma opinião pessoal de uma análise crítica, que aborde a complexidade do tema. 6.3.2.3 - Análise da Questão Discursiva 4 do Componente de Conhecimento Específico Na Tabela 6.19, constam as informações relativas à questão 4 do conjunto de questões do Componente de Conhecimento Específico.
O desempenho dos estudantes de todo o Brasil nesta questão foi o mais alto entre as três questões discursivas desse componente.
A Média geral do Brasil foi 32,4, sendo a maior Média registrada na região Sul ( 35,4 ), e a menor, na região Nordeste ( 29,1 ).
A nota Máxima 100,0 foi atingida em todas as regiões.
A Mediana em todo o Brasil foi 40,0.
Também foi zero a nota mínima de todas as regiões.
No Gráfico 6.12, está representada a distribuição de notas da questão discursiva 4, do Componente de Conhecimento Específico.
Essa distribuição tem moda na classe de estudantes no intervalo ( 30 ; 40 ], com frequência de 31,3 \%.
Outros intervalos que merecem destaque são ( 10 ; 20 ] e ( 50 ; 60 ], com frequências respectivamente iguais a 22,5 \% e 19,2 \%. 6.3.2.4 Comentários sobre as respostas à Questão Discursiva 4 Como na Questão 3, o enunciado da Questão 4 trazia um texto de opinião, publicado na mesma página sobre comunicação, marketing e mídia ( www.meioemensagem.com.br ), abordando marketing digital.
O comando solicitava que, a partir das ideias apresentadas na citação, o estudante redigisse um texto sobre o marketing digital abordando dois aspectos : mensuração dos dados atrelados ao comportamento do consumidor ; e segurança dos dados na rede, privacidade e conveniência.
A questão exigia tanto o conhecimento dos conteúdos e competências presentes nas Diretrizes Curriculares da área, como a capacidade crítica de articulá -los numa relação muito estreita entre mercado e o papel técnico e ético do publicitário.
A proposta de pedir uma análise do funcionamento do marketing digital, nos dias atuais, demandava uma análise sobre estratégias para alcançar o conhecimento, desenvolver atividades de comunicação integrada e monitorar e aferir resultados no mercado digital.
A proposta de redigir um texto possibilitava a análise da lógica estrutural do pensamento desenvolvido pelo estudante, mas, apesar da importância da temática, o fato de o comando ser bastante genérico abriu a possibilidade de diferentes interpretações por parte do estudante.
Os próprios conceitos a serem analisados ( como privacidade e comodidade ) são muito abertos.
Por conta da diversidade de possibilidades de reflexão sobre o tema e os conceitos envolvidos, grande parte das respostas não conseguiu relacioná -los de forma articulada às etapas da estratégia de marketing.
Pelo tema, muito presente nas discussões contemporâneas do marketing, a questão poderia ser considerada fácil.
No entanto, a subjetividade envolvida prejudicou a objetividade, a clareza e a organização das respostas.
De modo geral, os estudantes conseguiram atingir, ao menos, uma parte da reflexão esperada, aquela que diz respeito ao funcionamento do marketing digital e suas etapas, como a mensuração dos dados atrelados ao comportamento do consumidor.
No entanto, no que diz respeito aos outros itens ( como segurança de dados na rede, privacidade e conveniência ), o que se percebeu foi uma grande confusão em relação ao ponto de vista a ser adotado no desenvolvimento da argumentação.
Uma grande maioria respondeu à questão, a partir do ponto de vista do consumidor e não de um profissional da área de Publicidade e Propaganda, em vias de atuar no mercado de trabalho.
O mais comum foi os alunos abandonarem a análise prática e reflexiva da metodologia de funcionamento do marketing digital para o estudo do comportamento do consumidor e adotarem o ponto de vista do consumidor, cuja privacidade pode estar sendo invadida.
Em vez de pensar em estratégias no uso das tecnologias de comunicação, abordaram modos de defesa dessa estratégia.
Faltou também um maior aprofundamento das análises.
Por vezes, pontos importantes para a reflexão eram citados, mas não desenvolvidos.
Percebeu -se, também, uma dificuldade de elaboração de textos coerentes, com argumentação clara e concisa.
A resposta à Questão 4 foi deixada em branco por, aproximadamente, 12 \% dos participantes.
Houve uma expressiva quantidade de notas baixas, entre zero e 30, cerca de 33 \%.
No entanto, metade dos participantes ( 50 \% ) conseguiu atingir um patamar mediano, com notas entre 40 a 70.
Apenas uma minoria ( 5 \% ) conseguiu notas boas, acima de 75.
O desempenho demonstrado indica que o marketing digital, em todas as suas facetas e estratégias, pode ainda não estar integrado ao conteúdo programático do currículo dos cursos de Publicidade e Propaganda.
Os estudantes demonstraram dificuldade em tratar a questão de maneira crítica e aprofundar a discussão, para além de posições pessoais e do que estava presente no texto motivador. 6.3.2.5 Análise da Questão Discursiva 5 do Componente de Conhecimento Específico Na Tabela 6.20, constam as informações relativas à questão 5 do conjunto do Componente de Conhecimento Específico.
O desempenho dos estudantes de todo o Brasil nessa questão foi o mais baixo entre as três questões discursivas desse componente.
A nota Média dos estudantes de todo o Brasil foi 16,8.
A maior Média foi registrada na região Sudeste ( 18,3 ), enquanto a menor Média foi registrada na região Norte ( 10,7 ).
Quanto à variabilidade das notas, o Desvio padrão dos alunos do Brasil, como um todo, foi 18,9.
Enquanto o maior desvio foi encontrado na região Sudeste ( 19,5 ), o menor foi encontrado na região Norte ( 14,5 ).
A Mediana para o Brasil e para três das cinco regiões ( Sudeste, Sul e Centro-Oeste ) foi 10,0.
Para as demais regiões foi zero, indicando que, pelo menos, 50 \% dos participantes tiraram nota zero.
Também foi zero a nota mínima de todas as regiões.
Para o conjunto de alunos de Comunicação Social – Publicidade e Propaganda do Brasil, a nota Máxima foi 100,0, obtida nas regiões Nordeste, Sudeste e Sul.
As regiões Norte e Centro-Oeste, os alunos obtiveram nota Máxima 60,0.
No Gráfico 6.13, é representada a distribuição das notas da questão discursiva 5 do Componente de Conhecimento Específico.
Essa distribuição tem modal o intervalo correspondente aos que obtiveram nota zero, com frequência igual a 25,0 \% do total, que somada à frequência dos que deixaram essa questão em branco chega a 44 \% do total de participantes.
Outro intervalo que merece destaque é ( 20 ; 30 ], com 24,0 \%, constituindo -se em um máximo local. 6.3.2.6 Comentários sobre as respostas à Questão Discursiva 5 O enunciado da Questão 5 trazia um fragmento de texto de um portal da Internet sobre consumerismo e, a partir das ideias apresentadas no texto, solicitava que o estudante descrevesse “ dois ambientes estudados na análise macroambiental para compreender o comportamento de consumidores cada vez mais críticos e seletivos ”.
A questão era objetiva, e esperava -se do estudante uma resposta descritiva e bem específica.
A partir do conhecimento dos conceitos explicitados, conteúdo presente nas Diretrizes Curriculares da Área, a questão demandava a reflexão sobre as relações do marketing com o macroambiente, ou seja, com as variáveis incontornáveis que vêm de fora e interferem nas ações.
Desse modo, eram necessários conhecimentos não só relativos ao campo específico da propaganda, mas também à formação humana, político e social, parâmetros que são usados na Pesquisa de Mercado e estudos do Comportamento do Consumidor.
A questão demandava, apenas, que os conceitos fossem citados na descrição de dois ambientes, e não solicitava uma análise crítica.
Para responder a ela, adequadamente, era preciso conhecer o conceito de análise macroambiental.
Além disso, na escolha e descrição de dois ambientes, o texto introdutório sobre consumerismo induziu a escolha de exemplos relacionados com a questão do meio ambiente.
Tal situação acabou causando confusão em algumas respostas relativas à análise do ambiente macroambiental, que contemplaram a questão do meio ambiente ( ecologia ), esquecendo -se de que tal análise não se reduz a ele.
A grande maioria das respostas foi satisfatória, já que a questão solicitava a descrição de conceitos que fazem parte da grade curricular dos cursos.
No entanto, quando o conceito era desconhecido pelo estudante, foi difícil articular qualquer tipo de resposta.
Como já comentado, observou -se grande confusão causada pela contextualização da questão, que se referia ao meio ambiente.
Apesar de o meio ambiente ser uma das variáveis que podem ser analisadas na análise macroambiental, o padrão de resposta não previa a redução a esse ponto de vista.
No entanto, foi o que mais ocorreu.
Nesses casos, foi considerado como acerto o raciocínio de que o meio ambiente poderia ser uma variável externa a ser levada em conta nas pesquisas e estudos do Comportamento do Consumidor.
As respostas variaram pouco, basicamente em três grupos : as de estudantes que desconheciam totalmente o tema, e não se arriscaram a responder ou cometeram erros conceituais ; as daqueles que se concentraram na questão do meio ambiente, e, por fim, as daqueles que descreveram, de forma direta, um ou dois dos conceitos solicitados.
Os maiores erros encontrados dizem respeito ao não entendimento do que seria uma interferência dos ambientes externos nas práticas de consumo.
Nesses casos, mesmo quando citavam a questão do meio ambiente, que poderia ser uma interferência, não faziam essa associação.
Quase 20 \% dos estudantes participantes não responderam à questão.
Houve uma enorme quantidade ( 67 \% dos participantes ) de notas baixas, entre zero e 30.
Um número pequeno de respostas ( 13,4 \% ) conseguiu atingir um patamar mediano e recebeu notas entre 40 a 70.
Pouquíssimas respostas ( 0,6 \% ) atingiram notas boas, acima de 75.
As respostas revelam que muitos estudantes estão finalizando o curso sem conhecimento de conceitos básicos e que fazem parte dos currículos.
Ao mesmo tempo, mostram que estão bem informados acerca das discussões políticas e sociais, o que colabora para o pensamento sobre o próprio marketing.
No entanto, não conseguem desenvolver a reflexão sobre as informações que recebem cotidianamente.
Acabam por revelar que seria importante que as IES avaliassem a necessidade de atualização do que é ensinado, especialmente em maior sintonia com as práticas do mercado e, sempre, de modo a refletir sobre essas práticas. 6.3.3 Considerações Finais O desempenho dos estudantes avaliados no Enade-2018 pode ser considerado médio, a partir dos resultados da aplicação e correção da prova.
É possível notar que, mesmo quando demonstram conhecimento do tema, apresentam uma grande dificuldade na articulação de ideias, no desenvolvimento dos argumentos e das reflexões.
O incentivo a modos de estabelecer conexões entre as informações que atravessam seu cotidiano e a teoria poderia ser uma boa forma de desenvolver abordagens que levem a atitudes mais profissionais e menos baseadas em opiniões pessoais.
Chamou a atenção nas respostas em geral uma dificuldade de entender a partir de que ponto de vista a reflexão deveria ser realizada.
Muitas vezes, os estudantes confundiam o papel do publicitário, que cria estratégias de aproximação do consumidor com a marca, com o do consumidor, que se vê prejudicado em questões como a privacidade e a segurança.
Essa dificuldade prejudicou bastante o desenvolvimento dos argumentos nas questões discursivas de Conhecimentos Específicos.

% matched lemmas: aprendizado, centrado, destruir, esporádico, mentira, observação, personificar, peso, sócio-histórico-cultural, étnico-racial
\end{textsample}
